\documentclass[11pt]{article}
% -------------
% Preambule pro DRILY (cvičné úlohy ve stylu statnic) - MFF UK, UI / Strojové učení
% Sazba: pdflatex (UTF-8, čeština). Samostatné PDF na úroveň (1/2/3).
% -------------
\usepackage[T1]{fontenc}
\usepackage[utf8]{inputenc}
\usepackage{lmodern}
\usepackage[czech]{babel}
\usepackage[a4paper,margin=2.3cm]{geometry}

\usepackage{amsmath,amssymb,mathtools}
\usepackage{textcomp}
\usepackage[nopatch=all]{microtype}
\usepackage{xcolor}
\usepackage{enumitem}
\usepackage{titlesec}
\usepackage[most]{tcolorbox}
\usepackage{booktabs}
\usepackage{listings}
\usepackage{hyperref}

% - seznamy -
\setlist{leftmargin=1.4em, itemsep=2pt, topsep=3pt, parsep=0pt}

% - barvy -
\definecolor{mffblue}{RGB}{27,54,93}
\definecolor{mffaccent}{RGB}{0,90,156}
\definecolor{ulohagray}{RGB}{70,70,75}
\definecolor{resgreen}{RGB}{30,110,60}
\definecolor{gapred}{RGB}{150,30,30}
\definecolor{calorange}{RGB}{180,110,0}
\definecolor{codebg}{RGB}{245,245,245}

\titleformat{\section}{\Large\bfseries\color{mffblue}}{\thesection}{0.6em}{}
\titleformat{\subsection}{\large\bfseries\color{mffaccent}}{\thesubsection}{0.6em}{}
\titlespacing*{\section}{0pt}{1.4em}{0.6em}
\setcounter{secnumdepth}{2}
\setcounter{tocdepth}{1}

\hypersetup{colorlinks=true, linkcolor=mffaccent, urlcolor=mffaccent,
  pdftitle={Cvičné úlohy ke státnicím - Informatika / UI / Strojové učení},
  pdfauthor={Viktor Číhal}, bookmarksnumbered=true}

\renewcommand{\arraystretch}{1.2}
\setlength{\parindent}{0pt}
\setlength{\parskip}{5pt plus 2pt minus 1pt}
\emergencystretch=3em
\hbadness=10000

% - kód (C#); komentáře držíme bez diakritiky kvuli listings+utf8 -
\lstdefinestyle{cs}{
  basicstyle=\ttfamily\footnotesize, backgroundcolor=\color{codebg},
  keywordstyle=\color{mffaccent}\bfseries, commentstyle=\color{resgreen!70!black}\itshape,
  stringstyle=\color{gapred}, showstringspaces=false, columns=fullflexible,
  breaklines=true, frame=single, framerule=0pt, rulecolor=\color{codebg},
  xleftmargin=6pt, xrightmargin=6pt, aboveskip=4pt, belowskip=4pt,
  literate={ě}{{\v e}}1 {š}{{\v s}}1 {č}{{\v c}}1 {ř}{{\v r}}1 {ž}{{\v z}}1
           {á}{{\'a}}1 {í}{{\'\i}}1 {é}{{\'e}}1 {ý}{{\'y}}1 {ů}{{\r u}}1 {ú}{{\'u}}1,
  language=[Sharp]C}
\lstset{style=cs}

% =====================================================================
% Prostředí pro ÚLOHU a ŘEŠENÍ ve stylu statnic.
% =====================================================================

% čítač úloh
\newcounter{uloha}
\renewcommand{\theuloha}{\arabic{uloha}}

% \uloha{okruh}{téma} -> nadpis úlohy ve stylu zadání statnic
% okruh: "společná matematika" / "společná informatika" / "specializace UI-SU"
\newcommand{\ulohahead}[2]{%
  \refstepcounter{uloha}%
  \par\medskip\noindent
  {\large\bfseries\color{ulohagray}Úloha \theuloha.}\quad
  {\small\colorbox{ulohagray!12}{\textbf{#1}}}\quad{\bfseries #2}\par
  \vspace{2pt}}

% zadání: lehce rámovaný blok jako na zkoušce
\newtcolorbox{zadani}{enhanced, breakable, sharp corners=south,
  colback=ulohagray!5, colframe=ulohagray, boxrule=0pt, leftrule=3pt,
  left=8pt, right=8pt, top=4pt, bottom=5pt, before skip=4pt, after skip=4pt}

% podotázky se značkou bodu (kolik bodu daná část typicky nese)
\newcommand{\bod}[1]{\tikz[baseline=-0.55ex]\node[circle,fill=mffaccent,text=white,
  inner sep=1.3pt,font=\bfseries\scriptsize]{#1};}
% bez tikz (rychlejší): kroužek přes textový symbol
\renewcommand{\bod}[1]{{\footnotesize\colorbox{mffaccent}{\textcolor{white}{\bfseries #1\,b}}}}

% řešení
\newtcolorbox{reseni}{enhanced, breakable, colback=resgreen!6, colframe=resgreen,
  boxrule=0.6pt, arc=1.5pt, left=7pt, right=7pt, top=3pt, bottom=4pt,
  before skip=4pt, after skip=8pt, fonttitle=\bfseries\small, coltitle=resgreen,
  title={Řešení}}

% kalibrační poznámka: jak tato úloha sytí podlahu (common>=5 / spec>=7)
\newcommand{\kalibrace}[1]{\par\smallskip\noindent
  {\footnotesize\textcolor{calorange}{\textbf{Kalibrace:}}~\textit{#1}}\par\smallskip}

% past / častá chyba
\newcommand{\past}[1]{\par\smallskip\noindent
  \textbf{\color{gapred}Pozor (častá chyba):}~#1\par}

% info box
\newtcolorbox{infobox}[1][]{enhanced, breakable, colback=mffaccent!5, colframe=mffaccent,
  boxrule=0.7pt, left=6pt, right=6pt, top=4pt, bottom=4pt, arc=2pt,
  before skip=6pt, after skip=6pt, fonttitle=\bfseries, #1}


\begin{document}

\begin{center}
  {\LARGE\bfseries\color{mffblue} Programovací úlohy ke státnicím (C\#)}\\[4pt]
  {\small Informatika - Umělá inteligence, zaměření Strojové učení (UI-SU) \quad|\quad MFF UK}\\[2pt]
  {\footnotesize\itshape Implementační otázky společné informatiky přímo podle zadání zkoušek. Jazyk: C\# (volba majitele).}
\end{center}

\begin{infobox}[title={Proč zvlášť a jak to číst}]
\footnotesize
Ve společné informatice je skoro vždy jedna \textbf{implementační} otázka (návrh tříd nebo kód). To je nejčastější příčina \emph{nuly} na společné části - definici si vybavíš, ale ,,napište program'' tě zaskočí. Tenhle sešit sbírá opakované programovací typy ze zkoušek 2019-2026 s úplným řešením v C\#:
\begin{itemize}\setlength{\itemsep}{1pt}
  \item \textbf{Objektový návrh} (nejčastější): výrazový strom, maticová knihovna, graf - rozhraní, polymorfismus, princip open/closed, rozšiřitelnost.
  \item \textbf{Reprezentace dat}: čtení strukturovaného binárního souboru.
  \item \textbf{Vlákna a synchronizace}: producent-konzument semafory, oprava race condition.
  \item \textbf{Nízká úroveň}: alokátor paměti, ovladač zařízení (PIO), emulátor procesoru.
\end{itemize}
\textbf{Jak procvičovat:} přečti zadání, napiš kód \emph{na papír} (na zkoušce není počítač), pak porovnej. Cíl není odladěný program, ale \textbf{správná struktura} (rozhraní, klíčové metody, ošetření okrajových případu) - to nese body. Zelené \emph{Jádro} je kostra, kterou musíš umět z hlavy.\\
\textbf{Souvislost:} tyto úlohy odpovídají ,,no-zero'' drilum z úrovní 1-2; tady jsou pohromadě a s plnějším kódem.
\end{infobox}

\tableofcontents
\bigskip

\section{Objektový návrh}
\ulohahead{programování (C\#)}{Objektový návrh: výrazový strom (eval + symbolická derivace)}
\begin{zadani}
Navrhněte objektový model aritmetického výrazu nad konstantami, proměnnými a operacemi $+$ a $\cdot$. Umožněte (a) vyhodnocení v daném prostředí (hodnoty proměnných), (b) symbolickou derivaci podle proměnné, (c) rozšíření o novou operaci bez zásahu do stávajícího kódu. \emph{(Typ úlohy: 2019-jaro Q4, 2020-léto Q4 - výrazové stromy.)}
\end{zadani}

\begin{reseni}
Jedno rozhraní \texttt{IExpr}, každý druh uzlu je samostatná třída (polymorfismus, princip open/closed):
\begin{lstlisting}
public interface IExpr {
    double Eval(IReadOnlyDictionary<string,double> env);
    IExpr Derive(string x);          // symbolicka derivace podle x
}

public sealed class Const : IExpr {
    private readonly double v;
    public Const(double v) => this.v = v;
    public double Eval(IReadOnlyDictionary<string,double> env) => v;
    public IExpr Derive(string x) => new Const(0);
}

public sealed class Var : IExpr {
    private readonly string name;
    public Var(string name) => this.name = name;
    public double Eval(IReadOnlyDictionary<string,double> env) => env[name];
    public IExpr Derive(string x) => new Const(name == x ? 1 : 0);
}

public sealed class Add : IExpr {
    private readonly IExpr a, b;
    public Add(IExpr a, IExpr b) { this.a = a; this.b = b; }
    public double Eval(IReadOnlyDictionary<string,double> env) => a.Eval(env) + b.Eval(env);
    public IExpr Derive(string x) => new Add(a.Derive(x), b.Derive(x));
}

public sealed class Mul : IExpr {
    private readonly IExpr a, b;
    public Mul(IExpr a, IExpr b) { this.a = a; this.b = b; }
    public double Eval(IReadOnlyDictionary<string,double> env) => a.Eval(env) * b.Eval(env);
    // pravidlo o soucinu: (a*b)' = a'*b + a*b'
    public IExpr Derive(string x) => new Add(new Mul(a.Derive(x), b), new Mul(a, b.Derive(x)));
}
\end{lstlisting}
Použití: výraz $x\cdot x + 3$ je \texttt{new Add(new Mul(new Var("x"), new Var("x")), new Const(3))}; \texttt{Eval} s \texttt{\{["x"]=2\}} dá $7$, \texttt{Derive("x")} dá symbolicky $x{\cdot}1+1{\cdot}x$ ($+0$).

\textbf{Rozšiřitelnost (c):} novou operaci (např.\ \texttt{Sin}) přidáme jako novou třídu implementující \texttt{IExpr} - \emph{beze změny} stávajících tříd (open/closed). Sdílení podvýrazu (DAG) je zdarma: tentýž objekt \texttt{IExpr} lze odkázat z více míst (uzly jsou neměnné).
\end{reseni}

\begin{jadro}
Jedno rozhraní s operacemi (\texttt{Eval}, \texttt{Derive}), uzel = třída implementující ho; rekurze přes děti. Derivace součinu = pravidlo $(ab)'=a'b+ab'$. Rozšíření = nová třída, ne \texttt{switch}.
\end{jadro}

\past{Nepoužívej \texttt{switch} na typ uzlu - to porušuje open/closed a u nového uzlu zapomeneš větev. Polymorfismus přes rozhraní je správně. Uzly drž neměnné (immutable) kvuli sdílení.}

\ulohahead{programování (C\#)}{Knihovna pro práci s maticemi (rozhraní + hustá/řídká)}
\begin{zadani}
Navrhněte knihovnu pro matice: společné rozhraní, hustou implementaci a \emph{řídkou} (sparse) implementaci téhož rozhraní, a operaci násobení pracující nad rozhraním. \emph{(Typ úlohy: 2025-jaro Q2 - knihovna pro výpočty s maticemi.)}
\end{zadani}

\begin{reseni}
\begin{lstlisting}
public interface IMatrix {
    int Rows { get; }
    int Cols { get; }
    double this[int r, int c] { get; }   // indexer pro cteni prvku
}

public sealed class DenseMatrix : IMatrix {
    private readonly double[,] a;
    public DenseMatrix(double[,] a) => this.a = a;
    public int Rows => a.GetLength(0);
    public int Cols => a.GetLength(1);
    public double this[int r, int c] => a[r, c];
}

// ridka matice: ulozi jen nenulove prvky, ale navenek je to tataz abstrakce
public sealed class SparseMatrix : IMatrix {
    private readonly Dictionary<(int r, int c), double> nz = new();
    public int Rows { get; }
    public int Cols { get; }
    public SparseMatrix(int rows, int cols) { Rows = rows; Cols = cols; }
    public void Set(int r, int c, double v) {
        if (v != 0) nz[(r, c)] = v; else nz.Remove((r, c));
    }
    public double this[int r, int c] => nz.TryGetValue((r, c), out var v) ? v : 0.0;
}

// nasobeni pracuje jen pres rozhrani -> funguje pro libovolnou kombinaci
public static class MatrixOps {
    public static DenseMatrix Multiply(IMatrix x, IMatrix y) {
        if (x.Cols != y.Rows) throw new ArgumentException("nekompatibilni rozmery");
        var r = new double[x.Rows, y.Cols];
        for (int i = 0; i < x.Rows; i++)
            for (int j = 0; j < y.Cols; j++) {
                double s = 0;
                for (int k = 0; k < x.Cols; k++) s += x[i, k] * y[k, j];
                r[i, j] = s;
            }
        return new DenseMatrix(r);
    }
}
\end{lstlisting}
\textbf{Proč to takhle:} operace závisí jen na \texttt{IMatrix}, takže přidání nové reprezentace (pásová, jednotková, \dots) nevyžaduje změnu \texttt{Multiply} (polymorfismus přes indexer). Tovární metoda (\texttt{IMatrix Create(\dots)}) by volajícímu skryla výběr husté/řídké podle hustoty dat.
\end{reseni}

\begin{jadro}
Rozhraní \texttt{IMatrix} (rozměry + indexer); hustá i řídká jsou jen různé implementace; algoritmy (násobení) píšeš proti rozhraní, ne proti konkrétní třídě. Násobení $O(n^3)$, kontroluj rozměry.
\end{jadro}

\past{Násobení vyžaduje \texttt{x.Cols == y.Rows}. Řídká matice nesmí ukládat nuly (jinak ztrácí smysl); indexer vrací 0 pro chybějící klíč. Indexer je v C\# \texttt{this[int,int]}.}

\ulohahead{programování (C\#)}{Objektový návrh grafu a značkování komponent}
\begin{zadani}
Navrhněte rozhraní grafu a implementaci seznamem sousedu. Napište značkování komponent souvislosti (každému vrcholu přiřadit id komponenty). \emph{(Typ úlohy: 2022-podzim Q4, 2019-podzim Q2 - programování + grafy.)}
\end{zadani}

\begin{reseni}
\begin{lstlisting}
public interface IGraph {
    int N { get; }
    IEnumerable<int> Neighbors(int v);
}

public sealed class AdjListGraph : IGraph {
    private readonly List<int>[] adj;
    public AdjListGraph(int n) {
        adj = new List<int>[n];
        for (int i = 0; i < n; i++) adj[i] = new List<int>();
    }
    public int N => adj.Length;
    public void AddEdge(int u, int v) { adj[u].Add(v); adj[v].Add(u); } // neorientovany
    public IEnumerable<int> Neighbors(int v) => adj[v];
}

// znackovani komponent pomoci BFS; vraci comp[v] = id komponenty
public static int[] Components(IGraph g) {
    var comp = new int[g.N];
    Array.Fill(comp, -1);
    int c = 0;
    for (int s = 0; s < g.N; s++) {
        if (comp[s] != -1) continue;        // uz navstiveny
        var q = new Queue<int>();
        q.Enqueue(s); comp[s] = c;
        while (q.Count > 0) {
            int v = q.Dequeue();
            foreach (int w in g.Neighbors(v))
                if (comp[w] == -1) { comp[w] = c; q.Enqueue(w); }
        }
        c++;                                // dalsi komponenta
    }
    return comp;
}
\end{lstlisting}
\textbf{Složitost} $O(V+E)$ (každý vrchol/hrana zpracován jednou). Algoritmus závisí jen na \texttt{IGraph}, takže funguje i pro jinou reprezentaci (matice sousednosti) bez změny.
\end{reseni}

\begin{jadro}
Rozhraní \texttt{IGraph} (počet vrcholu + sousedé); reprezentace = seznam sousedu. Komponenty: pro každý nenavštívený vrchol spusť BFS/DFS a obarvi dosažené stejným id, $O(V+E)$.
\end{jadro}

\past{Algoritmus piš proti \texttt{IGraph}, ne proti konkrétní třídě. Nezapomeň na vrcholy v různých komponentách (vnější smyčka přes všechny vrcholy). Pole \texttt{comp} inicializuj na $-1$ (nenavštíveno).}


\section{Reprezentace dat a soubory}
\ulohahead{programování (C\#)}{Čtení strukturovaného binárního souboru}
\begin{zadani}
Binární soubor: 4 bajty ASCII ,,\texttt{REC1}'' (magic), \texttt{uint32} \texttt{count} (little-endian), poté \texttt{count} záznamů, každý \texttt{int32 id} + \texttt{float32 value} (LE). Napište načtení do pole struktur a ošetřete poškozený soubor. \emph{(Typ úlohy: 2024-podzim Q2, 2025-podzim Q2, reprezentace dat / binární soubor.)}
\end{zadani}

\begin{reseni}
\begin{lstlisting}
public readonly record struct Rec(int Id, float Value);

public static Rec[] Load(string path) {
    using var fs = File.OpenRead(path);
    using var br = new BinaryReader(fs);   // cte little-endian

    var magic = br.ReadBytes(4);           // ReadBytes muze vratit MENE bajtu
    if (magic.Length != 4 ||
        magic[0] != (byte)'R' || magic[1] != (byte)'E' ||
        magic[2] != (byte)'C' || magic[3] != (byte)'1')
        throw new InvalidDataException("spatny magic");

    uint count = br.ReadUInt32();
    // obrana proti nesmyslnemu count: over proti zbyvajici delce souboru
    long remaining = fs.Length - fs.Position;
    if ((long)count * 8 != remaining)
        throw new InvalidDataException("delka neodpovida count");

    var recs = new Rec[count];
    for (uint i = 0; i < count; i++) {
        int id = br.ReadInt32();           // pri zkraceni vyhodi EndOfStreamException
        float val = br.ReadSingle();
        recs[i] = new Rec(id, val);
    }
    return recs;
}
\end{lstlisting}
\textbf{Ruční složení int32 z bajtu} (princip endianity, nezávislé na CPU): \texttt{id = b0 | b1<<8 | b2<<16 | b3<<24}. Na big-endian CPU by selhalo až nativní přetypování (\texttt{BitConverter.ToInt32}), ne tohle explicitní složení; přenositelně lze použít \texttt{BinaryPrimitives.ReadInt32LittleEndian}.
\end{reseni}

\begin{jadro}
\texttt{BinaryReader} čte LE; ověř magic, \emph{ověř \texttt{count} proti délce souboru před alokací}, čti záznamy. \texttt{ReadBytes} vrací i méně bajtu (kontroluj délku), \texttt{ReadInt32} při zkrácení vyhodí výjimku.
\end{jadro}

\past{Nealokuj \texttt{new Rec[count]} bez ověření \texttt{count} (riziko OOM u poškozeného souboru). Little-endian = nejméně významný bajt první. Zarovnání ve \emph{formátu} souboru nepředpokládej.}


\section{Vlákna a synchronizace}
\ulohahead{programování (C\#)}{Vlákna: producent-konzument semafory a oprava race condition}
\begin{zadani}
(a) Napište producent-konzument s omezeným bufferem pomocí semaforu. (b) Ukažte časově závislou chybu (race condition) na sdíleném čítači a opravte ji. \emph{(Typ úlohy: 2024-léto Q2, 2026-jaro Q2 - Semafor; 2023-podzim Q1 - race condition.)}
\end{zadani}

\begin{reseni}
\textbf{(a) Producent-konzument} (buffer kapacity $N$): semafor \texttt{empty} počítá volná místa, \texttt{full} obsazená, \texttt{lock} chrání frontu.
\begin{lstlisting}
var empty = new SemaphoreSlim(N, N);   // volna mista
var full  = new SemaphoreSlim(0, N);   // obsazena mista
var gate  = new object();
var buffer = new Queue<int>();

void Produce(int item) {
    empty.Wait();                       // pockej na volne misto (P)
    lock (gate) { buffer.Enqueue(item); }
    full.Release();                     // pribyl prvek (V)
}
int Consume() {
    full.Wait();                        // pockej na prvek (P)
    int item; lock (gate) { item = buffer.Dequeue(); }
    empty.Release();                    // uvolnil se slot (V)
    return item;
}
\end{lstlisting}

\textbf{(b) Race condition:} \texttt{counter++} není atomické (load, inkrement, store); dvě vlákna se mohou proložit a jeden inkrement se ztratí. Oprava zámkem nebo atomickou operací:
\begin{lstlisting}
// chybne: for(...) counter++;  z N vlaken -> vysledek < N*kroku
// spravne A) zamek:
lock (gate) { counter++; }
// spravne B) atomicky:
Interlocked.Increment(ref counter);
\end{lstlisting}
\end{reseni}

\begin{jadro}
Semafor: \texttt{Wait}=P (sniž, blokuj při 0), \texttt{Release}=V. Pořadí: \emph{nejdřív} \texttt{empty.Wait()}/\texttt{full.Wait()}, \emph{pak} \texttt{lock}. Race: nechráněný \texttt{counter++} ztrácí inkrementy; oprava \texttt{lock} nebo \texttt{Interlocked}.
\end{jadro}

\past{Nezaměň pořadí semaforu a zámku (čekej na zdroj PŘED zámkem, jinak deadlock). \texttt{empty} a \texttt{full} se nesmí prohodit. \texttt{counter++} bez synchronizace je klasická race.}


\section{Nízkoúrovňové úlohy}
\ulohahead{programování (C\#)}{Jednoduchý správce paměti (first-fit alokátor)}
\begin{zadani}
Implementujte alokátor nad blokem o dané velikosti: \texttt{Alloc} strategií first-fit (s rozdělením bloku) a \texttt{Free} (se slučováním sousedních volných bloku). \emph{(Typ úlohy: 2025-léto Q2 - jednoduchý správce paměti.)}
\end{zadani}

\begin{reseni}
\begin{lstlisting}
public sealed class Allocator {
    public sealed class Block {           // slouzi i jako "handle" pro Free
        public int Offset, Size;
        public bool Free = true;
        public Block? Next, Prev;
    }
    private readonly Block head;
    public Allocator(int size) => head = new Block { Offset = 0, Size = size };
    private static int Align(int n, int a) => (n + a - 1) / a * a;

    public Block? Alloc(int n) {
        n = Align(n, 8);
        for (var b = head; b != null; b = b.Next)
            if (b.Free && b.Size >= n) {       // first-fit: prvni dost velky
                if (b.Size > n) Split(b, n);
                b.Free = false;
                return b;
            }
        return null;                           // nenalezeno
    }

    private void Split(Block b, int n) {       // odstrihni zbytek jako novy volny blok
        var rest = new Block { Offset = b.Offset + n, Size = b.Size - n,
                               Next = b.Next, Prev = b };
        if (b.Next != null) b.Next.Prev = rest;
        b.Next = rest; b.Size = n;
    }

    public void Free(Block b) {
        b.Free = true;
        if (b.Next is { Free: true } n) {      // sluc s pravym sousedem
            b.Size += n.Size; b.Next = n.Next;
            if (n.Next != null) n.Next.Prev = b;
        }
        if (b.Prev is { Free: true } p) {      // sluc s levym sousedem
            p.Size += b.Size; p.Next = b.Next;
            if (b.Next != null) b.Next.Prev = p;
        }
    }
}
\end{lstlisting}
\textbf{Fragmentace:} \emph{vnitřní} = zaokrouhlení na zarovnání nechá nevyužitý zbytek v bloku; \emph{vnější} = volná paměť roztříštěná na malé kousky. Slučování ve \texttt{Free} brání vnější fragmentaci.
\end{reseni}

\begin{jadro}
Spojový seznam bloku (offset, velikost, volný?). Alloc = první dost velký volný blok, případně rozděl. Free = označ volný a slučuj s volnými sousedy. Zarovnání zaokrouhlí velikost nahoru.
\end{jadro}

\past{Po \texttt{Free} se MUSÍ slučovat sousední volné bloky, jinak vnější fragmentace roste. Hlavička/režie bloku zabírá místo. First-fit je rychlý, ale ne nejúspornější (vs best-fit).}

\ulohahead{programování (C\#)}{Ovladač zařízení programovaným V/V (PIO)}
\begin{zadani}
Pro řadič disku s porty STATUS, COMMAND a DATA napište programem řízenou obsluhu (PIO): vydejte příkaz čtení a přečtěte 512 bajtu polling-em stavového registru. Vysvětlete variantu s přerušením. \emph{(Typ úlohy: 2024-jaro Q2 - ovladač pro řadič disku; 2021-podzim - myš.)}
\end{zadani}

\begin{reseni}
Řadič má registry adresované jako V/V porty; CPU s ním komunikuje instrukcemi in/out (zde abstraktně \texttt{Inb}/\texttt{Outb}).
\begin{lstlisting}
const int DATA = 0x1F0, STATUS = 0x1F7, COMMAND = 0x1F7;
const int BSY = 0x80, DRQ = 0x08;       // bity stavoveho registru

static byte Inb(int port) { /* HW: precti bajt z portu */ return 0; }
static void Outb(int port, byte v) { /* HW: zapis bajt na port */ }

static byte[] ReadSector() {            // PIO ctene polling-em
    Outb(COMMAND, 0x20);                // prikaz READ SECTOR
    var buf = new byte[512];
    for (int i = 0; i < 512; i++) {
        while ((Inb(STATUS) & BSY) != 0) { }   // cekej, az radic neni zaneprazdnen
        while ((Inb(STATUS) & DRQ) == 0) { }   // cekej, az jsou data pripravena
        buf[i] = Inb(DATA);                     // precti bajt dat
    }
    return buf;
}
\end{lstlisting}
\textbf{Polling vs přerušení:} polling \emph{aktivně} cyklí na stavovém registru a pálí CPU. \emph{Přerušení}: řadič po dokončení vyvolá IRQ; procesor přeruší běžící kód, skočí do obslužné rutiny (ISR), ta přečte data a oznámí dokončení (nastaví událost / probudí čekající vlákno). Procesor: dokončí instrukci, uloží kontext, podle vektoru přerušení skočí do ISR; OS: naplánuje obsluhu, po \texttt{iret} obnoví kontext.
\end{reseni}

\begin{jadro}
Registry řadiče = V/V porty (\texttt{in}/\texttt{out}). PIO polling: vydej příkaz, opakovaně čti STATUS (BSY/DRQ), pak čti DATA. Přerušení místo pollingu: IRQ $\to$ ISR přečte data $\to$ probudí čekajícího (nepálí CPU).
\end{jadro}

\past{Polling plýtvá CPU - u pomalých zařízení se použije přerušení. Pořadí bitu: čekej, až BSY=0 a DRQ=1, teprve pak čti DATA. Příkaz se zapisuje do COMMAND, data se čtou z DATA.}

\ulohahead{programování (C\#)}{Emulátor jednoduchého procesoru (instrukční dispatch)}
\begin{zadani}
Napište emulátor jednoduchého registrového procesoru s instrukcemi \texttt{LOAD}, \texttt{ADD}, \texttt{SUB}, \texttt{JMP}, \texttt{JZ}, \texttt{HALT}. Program je pole instrukcí, stav jsou registry a čítač instrukcí (PC). \emph{(Typ úlohy: 2022-léto Q4 - emulátor procesoru.)}
\end{zadani}

\begin{reseni}
\begin{lstlisting}
public enum Op { Load, Add, Sub, Jmp, Jz, Halt }
public readonly record struct Instr(Op Op, int A, int B);

public static int[] Run(Instr[] prog, int regCount) {
    var reg = new int[regCount];
    int pc = 0;
    while (true) {
        var ins = prog[pc];
        switch (ins.Op) {
            case Op.Load: reg[ins.A] = ins.B;          pc++;            break; // reg[A] = B
            case Op.Add:  reg[ins.A] += reg[ins.B];    pc++;            break; // reg[A]+=reg[B]
            case Op.Sub:  reg[ins.A] -= reg[ins.B];    pc++;            break;
            case Op.Jmp:  pc = ins.A;                                   break; // skok na adresu A
            case Op.Jz:   pc = (reg[ins.A] == 0) ? ins.B : pc + 1;      break; // if reg[A]==0 goto B
            case Op.Halt: return reg;
            default: throw new InvalidOperationException("neznama instrukce");
        }
    }
}
\end{lstlisting}
\textbf{Princip:} smyčka \emph{fetch-decode-execute} - načti instrukci na \texttt{pc}, rozskoč podle opkódu (\texttt{switch}), proveď, aktualizuj \texttt{pc} (skoky mění \texttt{pc} přímo). \emph{Alternativa OO:} každá instrukce jako třída implementující \texttt{IInstr.Execute(state)} - dispatch přes polymorfismus místo \texttt{switch} (snáz rozšiřitelné).
\end{reseni}

\begin{jadro}
Stav = registry + PC. Smyčka fetch-decode-execute: \texttt{switch} podle opkódu provede instrukci a posune PC; skoky nastaví PC přímo; \texttt{HALT} ukončí. OO varianta = instrukce jako třídy s \texttt{Execute}.
\end{jadro}

\past{Po skoku PC \emph{neinkrementuj} (skok ho už nastavil). \texttt{JZ} testuje registr a větví. Bez \texttt{HALT} / chybného PC hrozí nekonečná smyčka nebo přístup mimo pole - ošetři.}


\end{document}
